% Options for packages loaded elsewhere
\PassOptionsToPackage{unicode}{hyperref}
\PassOptionsToPackage{hyphens}{url}
\PassOptionsToPackage{dvipsnames,svgnames,x11names}{xcolor}
\documentclass[
  11pt,
]{article}
\usepackage{xcolor}
\usepackage[letterpaper,margin=0.8in]{geometry}
\usepackage{amsmath,amssymb}
\setcounter{secnumdepth}{-\maxdimen} % remove section numbering
\usepackage{iftex}
\ifPDFTeX
  \usepackage[T1]{fontenc}
  \usepackage[utf8]{inputenc}
  \usepackage{textcomp} % provide euro and other symbols
\else % if luatex or xetex
  \usepackage{unicode-math} % this also loads fontspec
  \defaultfontfeatures{Scale=MatchLowercase}
  \defaultfontfeatures[\rmfamily]{Ligatures=TeX,Scale=1}
\fi
\usepackage[]{lmodern}
\ifPDFTeX\else
  % xetex/luatex font selection
\fi
% Use upquote if available, for straight quotes in verbatim environments
\IfFileExists{upquote.sty}{\usepackage{upquote}}{}
\IfFileExists{microtype.sty}{% use microtype if available
  \usepackage[]{microtype}
  \UseMicrotypeSet[protrusion]{basicmath} % disable protrusion for tt fonts
}{}
\makeatletter
\@ifundefined{KOMAClassName}{% if non-KOMA class
  \IfFileExists{parskip.sty}{%
    \usepackage{parskip}
  }{% else
    \setlength{\parindent}{0pt}
    \setlength{\parskip}{6pt plus 2pt minus 1pt}}
}{% if KOMA class
  \KOMAoptions{parskip=half}}
\makeatother
\usepackage{longtable,booktabs,array}
\newcounter{none} % for unnumbered tables
\usepackage{calc} % for calculating minipage widths
% Correct order of tables after \paragraph or \subparagraph
\usepackage{etoolbox}
\makeatletter
\patchcmd\longtable{\par}{\if@noskipsec\mbox{}\fi\par}{}{}
\makeatother
% Allow footnotes in longtable head/foot
\IfFileExists{footnotehyper.sty}{\usepackage{footnotehyper}}{\usepackage{footnote}}
\makesavenoteenv{longtable}
\setlength{\emergencystretch}{3em} % prevent overfull lines
\providecommand{\tightlist}{%
  \setlength{\itemsep}{0pt}\setlength{\parskip}{0pt}}
\usepackage{fancyhdr,enumitem,xurl,microtype,needspace,etoolbox}
\setlength{\parindent}{0pt}
\setlength{\parskip}{5pt plus 1pt}
\setlength{\headheight}{14pt}
\setlength{\emergencystretch}{2em}
\clubpenalty=10000
\widowpenalty=10000
\setlist{leftmargin=1.4em,itemsep=2pt,topsep=3pt}
\pretocmd{\subsection}{\Needspace{7\baselineskip}}{}{}
\pretocmd{\subsubsection}{\Needspace{12\baselineskip}}{}{}
\pagestyle{plain}
\AtBeginDocument{\hypersetup{pdftitle={Massive QCD: verification assessment},
pdfauthor={Review prepared with Codex},linkcolor=blue!45!black,urlcolor=blue!45!black}}
\usepackage{bookmark}
\IfFileExists{xurl.sty}{\usepackage{xurl}}{} % add URL line breaks if available
\urlstyle{same}
\hypersetup{
  colorlinks=true,
  linkcolor={Maroon},
  filecolor={Maroon},
  citecolor={Blue},
  urlcolor={Blue},
  pdfcreator={LaTeX via pandoc}}

\author{}
\date{}

\begin{document}

\begin{center}
{\Large\bfseries Massive QCD: verification assessment}\par
{\normalsize Summary of results and checks\quad\textbullet\quad20 November 2026}
\end{center}

\textbf{Results.} The submitted proof constructs a continuum field
system for four-dimensional SU(3) quantum chromodynamics (QCD) with six
massive quark flavors. It is written in Lean, a computer proof
assistant, and passed verification of its precise mathematical
statement. It builds on the framework used in OpenAI's November 4 (Nov4)
pure Yang-Mills submission, which treats gauge fields without quarks.
The QCD construction supplies continuum correlations, six positive
distinct fitted mass parameters, identities for currents and the stress
tensor, and a reconstruction of real-time quantum fields, all within one
construction.

Our additional checks establish explicit lattice-to-continuum
correlation limits at distinct points, confirm that the properties hold
for the same constructed objects, and bound normalization changes. They
also confirm the expected one-loop dependence on the number of quark
flavors in the source calculations. We found that the formal
mass-fitting rule and the rule described in the separate informal review
can differ on synthetic data; this does not establish a discrepancy in
the constructed theory.

\textbf{Checks supplied by OpenAI and rerun here.}

\begin{itemize}
\tightlist
\item
  Check source files against recorded hashes, prepare fixed dependency
  versions, and build the complete proof.
\item
  Use Comparator to check the theorem statement and its definitions
  against a preserved reference, and check that the reference has not
  changed.
\item
  Check the allowed axioms and verify the proof with two proof checkers:
  Nanoda and Lean's kernel.
\end{itemize}

\textbf{Additional checks developed for this review.}

\begin{itemize}
\tightlist
\item
  \textbf{Lean proofs:} lattice correlation limits, properties of one
  common construction and its selected vacuum, current identities,
  normalization bounds, and a conditional test of the stress-tensor
  criterion.
\item
  \textbf{Mass-fit algebra:} the exact solution and its stability,
  comparison with the informal prescription, and 44 exact-arithmetic
  diagnostics.
\item
  \textbf{Lattice and flavor checks:} 41 numerical checks of the Wilson
  quark operator; Lean checks of one-loop flavor dependence and a
  fermion-determinant calculation in a fixed background.
\item
  \textbf{Failure detection:} deliberately altered inputs, statements,
  definitions, axioms and coefficients, together with checks of the
  stored evidence.
\end{itemize}

All completed checks passed within their stated scopes. Numerical tests
are separate from formal proofs; proofs that use the submitted theorem
are not a second construction of QCD.

\textbf{Scope and open questions.} Full identification with the
conventional physical operators, currents, masses and coupling remains
open, as does comparison with a complete free gauge-and-quark model.
Arbitrary prescribed quark masses, particle masses, a mass gap and
confinement are not established. \textbf{Derivation of the physical
two-loop coefficient is deferred.} The run used macOS without Linux
process isolation. These formal checks do not verify every step of the
separate informal argument.

\clearpage

\section{Detailed assessment}\label{detailed-assessment}

This assessment concerns {\small \nolinkurl{massive-qcd-lean}}, the
supplied collection of Lean proof files and verification tools. It
reports the mathematical checks completed on 17-18 November 2026 and a
source comparison made on 19 November. The description of the regulator
and construction was expanded on 20 November. It was prepared with Codex
and is not an external peer-review certification. The submitted source
and theorem statement have not been changed.

Three bodies of work are distinguished throughout: \textbf{the submitted
formal proof}, \textbf{the additional checks written for this review},
and \textbf{a separate informal proof review}. The last is an analytic
review of the prose argument; it did not examine the QCD Lean code. A
result from one of these reviews does not automatically settle the
questions addressed by the others.

\subsection{Overall judgment}\label{overall-judgment}

\textbf{The received proof passed the supplied formal acceptance
procedure.} This is convincing evidence that it proves the specified
Lean theorem, subject to the recorded software and runtime assumptions.
The check compared the theorem and its definitions with a preserved
reference, restricted the permitted axioms, and obtained acceptance from
both Nanoda and Lean's kernel. A kernel is the component that checks the
logical steps of a formal proof. This provides more evidence than
compiling the source alone.

There is a further question: does the precise formal statement capture
the intended physical claim about QCD? The theorem has substantial
content. It requires a lattice construction, nonzero continuum fields,
six positive distinct mass parameters, stress-tensor and current
identities, and real-time reconstruction. These properties must belong
to the same construction. The final theorem does not simply take the
desired renormalization-group estimates or reconstruction conclusions as
assumptions.

Our added checks clarify several of these requirements and their
consequences. They also test independent algebra, finite lattice
operators and the verifier's ability to reject deliberate errors. All
completed checks passed. A complete independent assessment of the formal
definitions' physical meaning, and of every estimate in the informal
argument, remains unfinished.

\subsection{Relationship to the Nov4 pure Yang-Mills
construction}\label{relationship-to-the-nov4-pure-yang-mills-construction}

Pure Yang-Mills theory describes gauge fields without quarks. QCD
couples the SU(3) gauge field to quark fields. The QCD development
\textbf{builds on the pure Yang-Mills framework also used in the Nov4
submission}: it reuses lattice definitions and substantial analytic and
construction machinery. The QCD-specific work introduces the lattice
quark operator and fermionic variables, flavor and mass dependence,
current identities, and the estimates needed for the coupled
gauge-and-quark construction.

The
\href{https://github.com/yang-mills-group/QCD/blob/main/evidence/source-relation/nov4-overlap.json}{source
comparison} checked the archived QCD files against the Nov4 release
manifest, which lists file hashes. Of the QCD package's 7,594 Lean
modules, \textbf{3,667 have matching relative paths in Nov4 and 3,018
are byte-identical}; 649 shared paths have different contents. For
example, {\small \nolinkurl{YM/Wilson/Law.lean}} is unchanged. All 3,018
identical modules are reached by following imports from the QCD
theorem's entry file, {\small \nolinkurl{YM/Plata/Massive.lean}}. This
establishes shared source infrastructure. Imports can include
declarations not used by the final proof, and source overlap alone does
not establish when or in which direction files were copied.

The QCD proof has its own final theorem,
{\small \nolinkurl{YM.Massive.interacting_six}\ \nolinkurl{:}\ \nolinkurl{YM.Massive.Goal}},
and its own successful verification. It does not directly invoke Nov4's
final pure-YM theorem, {\small \nolinkurl{YangMills.jointExistence}}, or
its perturbative-matching theorem,
{\small \nolinkurl{YangMills.perturbativeMatching}}. In particular, the
earlier submission's conclusions about agreement with perturbation
theory to all orders do not automatically become QCD conclusions. The
shared machinery is reused, but the QCD statement and proof must be
assessed on their own terms.

\subsection{What the QCD theorem says}\label{what-the-qcd-theorem-says}

The final theorem is named
{\small \nolinkurl{YM.Massive.interacting_six}}; its required statement
is named {\small \nolinkurl{YM.Massive.Goal}}. The statement says that,
for every positive tolerance, all sufficiently large values of the
reference parameter \(\beta\) admit a construction satisfying
{\small \nolinkurl{TheoryAt}}. The tolerance is fixed before the
threshold for that parameter is chosen
({\small \nolinkurl{YM/Massive/Target/Statement.lean:39-58}}). It is an
existence result, not a computation of the measured quark masses.

{\small \nolinkurl{TheoryAt}} collects the required properties of
\textbf{one common construction}. In Lean, the objects supplied by an
existence theorem are often called its \emph{witnesses}. Here they
include the lattice sequence, field normalization, continuum fields,
mass readings, stress tensor, correlation functions and vacuum family.
The statement requires:

\begin{itemize}
\tightlist
\item
  a common sequence of lattice spacings, volumes, couplings and bare
  parameters;
\item
  field normalization and mass readings along that sequence;
\item
  continuum limits of the same lattice correlations, with the required
  Euclidean properties;
\item
  axial and vector current identities, with six positive distinct fitted
  masses;
\item
  stress-tensor, symmetry and field-content conditions;
\item
  reconstruction of real-time fields and vacuum sectors from those same
  data.
\end{itemize}

Euclidean correlations use imaginary time. Reconstruction produces the
Hilbert space and real-time quantum fields from those correlations. The
requirement that all properties use the same objects matters: separate
existence statements could otherwise refer to different theories.

The setting is four dimensions, gauge group SU(3) in its defining
representation, and six fundamental Dirac flavors. A flavor is a quark
species. The masses are \textbf{scheme masses}: parameters determined by
the construction's particular measurement and fitting prescription. They
are not automatically pole masses or hadron masses. The target does not
allow arbitrary prescribed quark masses as inputs, and explicitly does
not request a particle spectrum, a spectral gap, confinement or
color-charged fields ({\small \nolinkurl{Statement.lean:13-29}}). It
permits changes to the renormalization-group estimate scheme used in the
informal argument. A comment marking the reference target as
``UNPROVED'' identifies the problem to be solved; the solution is in the
submitted proof files.

The stress tensor encodes responses to spacetime symmetries. The target
requires a nontrivial stress tensor, a uniqueness condition on a
specified class of conserved candidates, and nonzero translation and
rotation actions ({\small \nolinkurl{Physics.lean:137-155}}). The
reconstructed fields must have a normalized cyclic vacuum, covariance,
forward energy-momentum spectrum, a unique translation-invariant line, a
faithful representation and graded spacelike commutation
({\small \nolinkurl{Wightman.lean:80-116}}). These are substantial
mathematical requirements. Their full comparison with the conventional
physical definitions is a separate part of the review.

\subsection{What the quark extension
adds}\label{what-the-quark-extension-adds}

Adding quarks changes the full construction, beyond their contribution
to the running coupling. The main additional obligations are:

\begin{itemize}
\tightlist
\item
  \textbf{Fermionic integration and renormalization-group blocking.} The
  cutoff theory includes anticommuting, or \emph{Grassmann}, quark
  variables and their integration rule. Exact blocking combines fine
  lattice degrees of freedom into coarser ones while retaining the
  fermion determinants and field insertions. For a Gaussian fermion
  density at a fixed gauge background, eliminating fine fermion
  variables produces a Schur complement (an effective Dirac matrix) and
  a determinant prefactor. The coupled construction must control fermion
  propagation, determinant contributions and generated interactions as
  the cutoff is removed. Its full coarse density is not restricted to
  the original Wilson action with changed parameters.
\item
  \textbf{Masses, currents and local-field normalization.} The Wilson
  term breaks chiral symmetry at nonzero lattice spacing, so mass tuning
  includes an additive counterterm. The construction must normalize and
  mix quark operators and establish the continuum axial and vector
  identities used to define its mass parameters. Equality of bare and
  renormalized masses is not assumed.
\item
  \textbf{Reflection positivity and partition normalization.} Reflection
  positivity, which underlies a positive Hilbert-space inner product,
  must be extended to the gauge-and-fermion theory. The supplied
  finite-lattice result covers gauge-invariant insertions, including
  those with an odd number of fermionic factors, and reflections across
  antiperiodic seams, under its lattice-size and coupling conditions
  ({\small \nolinkurl{YM/Massive/Finite/Squares.lean}}). Strict
  positivity of the partition function is a separate requirement: for
  unequal masses, reality of the fermion determinant product alone does
  not prove a pointwise nonnegative gauge weight.
\item
  \textbf{The enlarged continuum theory.} Quark bilinears, currents and
  baryon operators join the gauge observables. Their jointly normalized
  correlations must have continuum limits, and reconstruction must
  respect fermionic signs. The stress tensor includes matter
  contributions, and the symmetry, mass and stress assertions must
  survive in the same selected vacuum theory. These requirements extend
  the pure-YM observable and reconstruction machinery.
\end{itemize}

This is a description of the submitted construction and its
requirements, read from the source; it is not a list of new independent
verification results. The
\href{https://github.com/yang-mills-group/QCD/blob/main/provenance/supplied-comparator/frozen/reference/YM/Massive/Target/Tower.lean}{RG
target},
\href{https://github.com/yang-mills-group/QCD/blob/main/provenance/supplied-comparator/frozen/reference/YM/Massive/Target/Calibration.lean}{observable
calibration} and
\href{https://github.com/yang-mills-group/QCD/blob/main/provenance/supplied-comparator/frozen/reference/YM/Massive/Target/Physics.lean}{mass/current
target} specify the relevant objects. Our additional checks and their
limits are reported below.

\subsubsection{The cutoff action and its
parameters}\label{the-cutoff-action-and-its-parameters}

The regulator uses \textbf{nearest-neighbor Wilson Dirac fermions with
Wilson parameter \(r=1\)}. Each of the six flavors has four spin
components and three colors in the fundamental representation of SU(3).
In dimensionless lattice units, its action is

\[
S=\beta\sum_p\left(1-\frac13\operatorname{Re}\operatorname{Tr}U_p\right)
  +\sum_{f=1}^{6}\bar\psi_f D_W(U,m_f^{\mathrm{bare}})\psi_f,
\]

where \(U_p\) is the ordered product of links around a plaquette, and
the barred and unbarred quark variables are independent Grassmann
generators. Expanding the
\href{https://github.com/yang-mills-group/QCD/blob/main/provenance/supplied-comparator/frozen/reference/YM/QCD/Lattice/Stencil.lean\#L129}{source's
Wilson stencil} gives

\[
D_W(U,m)=(m+4)I-\frac12\sum_{\mu=1}^{4}
\left[(1-\gamma_\mu)T_\mu(U)+(1+\gamma_\mu)T_\mu(U)^\dagger\right].
\]

Here \(\gamma_\mu\) are Euclidean spin matrices and \(T_\mu(U)\) is the
forward covariant shift, including the boundary sign. The normalization
of the gauge term follows
\href{https://github.com/yang-mills-group/QCD/blob/main/provenance/supplied-comparator/frozen/reference/YM/Wilson/Law.lean\#L29}{the
Wilson energy} and
\href{https://github.com/yang-mills-group/QCD/blob/main/provenance/supplied-comparator/frozen/reference/YM/Gauge/MatrixGroup.lean\#L119}{the
plaquette price}.

At a fixed lattice size, the
\href{https://github.com/yang-mills-group/QCD/blob/main/provenance/supplied-comparator/frozen/reference/YM/Massive/Finite/Core.lean\#L27}{regulator
parameters} are \(\beta\), one real bare mass per flavor, and the
positions of the boundary seams. The flavor mass matrix is diagonal.
\textbf{Fermions are antiperiodic in all four directions; gauge links
are periodic.} The actual cutoff sequence places the seams at coordinate
zero. Bare masses may have either sign. Setting \(m=0\) gives the
tree-level massless stencil; a renormalized massless theory would still
require tuning. The final six positive distinct fitted masses are
outputs of the continuum current prescription, not arbitrary prescribed
masses or copies of the bare parameters.

The joint density is \(e^{-S}\). Integrating its Grassmann variables
gives the gauge weight

\[
e^{-S_g(U)}\prod_{f=1}^{6}\det D_W(U,m_f^{\mathrm{bare}}),
\]

where \(S_g\) is the plaquette term above. Thus the quark determinants
are retained. Along each RG run, the source also specifies a
scale-dependent \(\beta\), a positive spinor rescaling \(w\), and the
full blocked density ({\small \nolinkurl{Run}} and
{\small \nolinkurl{Run.Actual}}). Its comparison of the running coupling
with a measured stiffness is discussed below.

\subsubsection{Gauge links and smooth test
connections}\label{gauge-links-and-smooth-test-connections}

The regulated action and exact blocking act directly on lattice SU(3)
links. Gauge transport uses products along discrete paths, and the
fermion block averages include those transports. These operations do not
require assigning a continuous spacetime gauge potential to each lattice
configuration being integrated. See the
\href{https://github.com/yang-mills-group/QCD/blob/main/provenance/supplied-comparator/frozen/reference/YM/RG/Block.lean\#L31}{gauge
blocking map} and
\href{https://github.com/yang-mills-group/QCD/blob/main/provenance/supplied-comparator/frozen/reference/YM/QCD/Flume/Collect.lean\#L39}{fermion
block average}.

\textbf{Smooth auxiliary connections are used for observable
calibration.} From prescribed local Taylor data, called \emph{jets}, the
source constructs a polynomial gauge potential and solves its
parallel-transport differential equation along links. The massive
calibration also uses a smooth cutoff that equals one near the chosen
origin and vanishes outside a larger neighborhood. It evaluates lattice
insertions, including transported insertions, on the resulting test
links and matches Taylor coefficients to the desired continuum operator
labels. This is the direction \textbf{smooth test connection to lattice
holonomies}, rather than an interpolation of every fluctuating lattice
configuration. The definitions are
\href{https://github.com/yang-mills-group/QCD/blob/main/provenance/supplied-comparator/frozen/reference/YM/Source/Local.lean\#L24}{the
polynomial connection and holonomy} and
\href{https://github.com/yang-mills-group/QCD/blob/main/provenance/supplied-comparator/frozen/reference/YM/Massive/Target/Calibration.lean\#L67}{the
smooth cutoff and calibration test}.

Likewise, an insertion being \emph{continuous in the gauge links} refers
to dependence on the finite collection of link matrices, not to
spacetime continuity of a gauge potential. The continuum output is
obtained through correlation limits and quantum-field reconstruction.
The target requires uniform convergence of normalized lattice
correlations away from coincident insertion points
(\href{https://github.com/yang-mills-group/QCD/blob/main/provenance/supplied-comparator/frozen/reference/YM/Massive/Target/Euclidean.lean\#L140}{{\small \nolinkurl{Fields.Scaling}}});
it does not define that output by interpolating sampled gauge fields
into continuum connections. Matching the calibrated continuum fields
fully to conventional physical operators remains part of the open
interpretation work.

\subsection{What OpenAI's supplied checks
establish}\label{what-openais-supplied-checks-establish}

The package includes a verification workflow as well as the mathematical
proof. We ran that workflow. Its main steps are:

\begin{enumerate}
\def\labelenumi{\arabic{enumi}.}
\tightlist
\item
  \textbf{Preserve the problem and inputs.} File hashes, or digital
  fingerprints, check consistency with the recorded source. Fixed
  dependency versions and a separate reference build control which
  definitions are used.
\item
  \textbf{Compare the solution with the reference.} The \emph{frozen
  target} is the preserved theorem statement and the definitions it
  depends on. Comparator checks that the submitted solution proves that
  target, rather than a changed statement or a changed definition with
  the same name. The reference export must also match its expected hash.
\item
  \textbf{Check the proof and its assumptions.} The workflow restricts
  the allowed axioms and checks the proof with Nanoda and Lean's kernel.
  This guards against a new axiom or an admitted proof gap being
  accepted as the requested solution.
\end{enumerate}

We found no separate supplied numerical physics test suite. Some source
files named {\small \nolinkurl{Checker}} concern mathematical
checkerboard constructions; {\small \nolinkurl{YM/QCD/Port/Bench.lean}}
is mathematical source, not a record of an independent physics
benchmark. The submitted proofs already include Clifford and projector
identities, gauge covariance and gamma-five hermiticity. Our
finite-operator tests therefore use an independent implementation rather
than merely restating those lemmas.

These checks are strong evidence for formal correctness of the specified
theorem. They do not decide whether a definition named ``mass'' or
``coupling'' is the physical quantity a reader expects. That requires
examining the definitions and their relation to observables. Hash checks
also establish consistency of source bytes, not independently
authenticated authorship.

\subsection{Additional checks developed for this
review}\label{additional-checks-developed-for-this-review}

The following tests were written for this review and were not part of
the received package. Some prove more explicit consequences of the
submitted theorem; others independently verify algebra, lattice
conventions or error detection. A consequence proof can reveal a
specification problem or establish a useful implication without
providing a second independent construction of QCD.

\subsubsection{Keeping the properties on one
construction}\label{keeping-the-properties-on-one-construction}

An existence theorem must supply compatible objects, not just separate
examples with different desired properties. Our
{\small \nolinkurl{AcceptanceSpec.lean}} restates every clause of
{\small \nolinkurl{TheoryAt}} using explicit common objects, then adds a
selected vacuum reconstruction for those objects.
{\small \nolinkurl{AcceptanceProof.lean}} derives this statement from
the submitted theorem.

The selected vacuum sector retains the required Euclidean, symmetry and
comparison properties. The masses remain positive and distinct. The
off-diagonal pseudoscalar fields remain nonzero, as do their
reconstructed operators; faithfulness of the representation ensures that
a nonzero field is not lost in reconstruction. Separate lemmas check
that the mass sums and differences used as denominators in the
normalization formulas are nonzero.

All four theorems compiled with only the permitted standard axioms. They
check the extraction of a single compatible construction from the
submitted result. They do not independently rebuild the theory.

\subsubsection{The mass fit and the difference from the informal
prescription}\label{the-mass-fit-and-the-difference-from-the-informal-prescription}

The six mass parameters are obtained by fitting pairwise readings. Write
\(b_{ij}\) for the reading associated with flavors \(i\) and \(j\), and
\(u_i\) for a candidate fitted mass. The formal definition minimizes

\[
 L_b(u)=\sum_{i\ne j}(b_{ij}-u_i-u_j)^2,
\]

using all 30 ordered off-diagonal pairs: six choices of \(i\) and five
different choices of \(j\) ({\small \nolinkurl{Physics.lean:67-87}}).
Our independently stated solution in
\href{https://github.com/yang-mills-group/QCD/blob/main/checks/MassFitSpec.lean}{MassFitSpec}
is

\[
 m_i^*=\frac18\sum_{j\ne i}(b_{ij}+b_{ji})-
          \frac1{40}\sum_{k\ne j}b_{kj}.
\]

The new Lean proof establishes, for arbitrary real readings, including
\(b_{ij}\ne b_{ji}\),

\[
 L_b(u)=L_b(m^*)+8\sum_i(u_i-m_i^*)^2+
                   2\left(\sum_i(u_i-m_i^*)\right)^2.
\]

The added terms are nonnegative and vanish together only at \(u=m^*\).
Thus the formula gives the unique global minimizer of the \textbf{actual
submitted loss function}, without using a mass-fit theorem from the
submission. Four mass-fit theorems, including a strict comparison
example and its error bound, compiled with the permitted standard
axioms.

A separate exact-rational calculation supplies 44 diagnostics. It checks
the least-squares matrix identity \(A^{\mathsf T}A=8I+2J\), where \(J\)
is the all-ones matrix; the normal equations on all 36 input basis
arrays; the irrelevance of diagonal readings; and detection of an
intentionally wrong divisor. It also checks stability: reading errors
bounded by \(\delta\) change each fitted mass by at most \(3\delta/2\)
and each mass difference by at most \(2\delta\). At the target's error
allowance, \(\delta=\mathrm{margin}/100\), the remaining positivity and
separation margins are at least 0.985 and 0.98 times their reference
margins. These calculations check the fitting rule, not the construction
of the continuum readings themselves.

The separate informal review, revision 0.2, page 8, describes a
different rule using only the six unordered pairs
\(\{12,13,23,14,15,16\}\). A concrete example distinguishes the rules.
Start with masses \(r=(1,2,3,4,5,6)\) and readings \(b_{ij}=r_i+r_j\),
then add \(\delta=1/1000\) to both orders of pair \(\{5,6\}\). The six
selected pairs do not see this change and still return \(r\). The
all-pairs fit subtracts \(\delta/20\) from the first four masses and
adds \(\delta/5\) to the last two. Its loss is \((6/5)\delta^2\), below
the six-pair answer's \(2\delta^2\).

\textbf{The two prescriptions are therefore not identical on all
sufficiently small perturbations.} We have not shown that the actual
constructed readings exhibit this difference. The finding identifies a
difference between prescriptions that should be documented; it does not
refute the formal theorem.

The target also uses pair-dependent rescalings in its axial-current
relations, often called \emph{PCAC} relations (partial conservation of
the axial current). Once the underlying identities and nonzero
denominators are established, these rescalings produce the specified
mass sums and differences ({\small \nolinkurl{Physics.lean:94-117}}).
This alone does not identify the resulting density normalization with a
conventional unrescaled Wilson or pion operator, or identify a particle
mass.

\subsubsection{Independent checks of the finite lattice quark
operator}\label{independent-checks-of-the-finite-lattice-quark-operator}

The Wilson quark operator specifies fermion propagation on the regulated
lattice. Our Python program,
{\small \nolinkurl{scripts/test_wilson.py}}, independently implements
the spin matrices, antiperiodic boundary condition and Wilson lattice
operator specified in {\small \nolinkurl{YM/QCD/Lattice/Gamma.lean}} and
{\small \nolinkurl{Stencil.lean}}. It compares the position-space
operator with a separately written momentum-space formula for the free
operator:

\[
 D(p)=\left(m+\sum_a(1-\cos p_a)\right)I+i\sum_a\gamma_a\sin p_a.
\]

Here \(p\) is lattice momentum, \(m\) is the lattice mass parameter, and
\(\gamma_a\) are the Euclidean spin matrices. The 41 checks cover their
Clifford algebra and chirality conventions, the allowed antiperiodic
momenta on lattices of side length 2 and 4, two masses, gauge covariance
for noncommuting determinant-one links, gamma-five hermiticity, and the
unequal-mass axial identity including the Wilson term. Deliberately
changing the boundary sign, swapping projectors or omitting the Wilson
term causes the intended failures.

These are numerical diagnostics with specified tolerances and recorded
source hashes. The correspondence between the Python implementation and
Lean source was reviewed manually; there is no formally proved
translation between them. The tests do not establish interacting
continuum limits, reflection positivity of the full measure,
renormalization-group estimates, or disappearance of the remaining
Wilson correction in a continuum identity.

\subsubsection{Checking that the verifier detects
errors}\label{checking-that-the-verifier-detects-errors}

Nine input-guard controls accept the recorded preparation and reject
altered proof or reference files, changed manifests or solution files,
an extra source file, a source symlink, and a changed staged reference.
Five Comparator controls accept a correct theorem and reject a changed
statement, an extra axiom, Lean's proof-gap axiom
{\small \nolinkurl{sorryAx}}, and a changed definition body. A rejection
must include the expected diagnostic; an unrelated crash is not counted
as a successful test.

The first three review targets also passed Comparator and both kernels:
correctness of the mass fit, the strict comparison with the six-pair
prescription, and extraction of one common construction. Their status is
recorded in
\href{https://github.com/yang-mills-group/QCD/blob/main/evidence/independent-checks/review-comparator/result.json}{the
review receipt}. These new review modules were compiled using the
previously built submission libraries; this was not a second clean build
of the full submission.

\subsection{Further checks connecting the definitions to physical
requirements}\label{further-checks-connecting-the-definitions-to-physical-requirements}

The
\href{https://github.com/yang-mills-group/QCD/blob/main/docs/semantic-tests.md}{four-area
correspondence suite} adds 33 audited theorems and seven Comparator
targets. ``Correspondence'' here means checking how the submitted
definitions relate to a more explicit requirement, such as a lattice
correlation limit or a test intended to distinguish an interacting
theory from a free one.

\textbf{Lattice correlation limits.} The proofs derive convergence of
the actual normalized lattice correlations when lattice sites approach
distinct continuum points. The same choice of approximating sites works
for all allowed field insertions at that configuration. This uses the
same lattice sequence, calibration and correlation function as the
accepted theorem. It does not establish limits when points collide, full
distributional limits, or a direct lattice realization of each
individually selected vacuum sector.

\textbf{A test of the stress-tensor criterion.} Under the target's exact
candidate and uniqueness definitions
({\small \nolinkurl{Physics.Candidate}} and
{\small \nolinkurl{Stress.unique}}), two independently detectable
conserved sectors cannot satisfy the required uniqueness condition. A
concrete degenerate model is also formally rejected. To turn this into a
complete free-QCD exclusion test, one would still have to construct the
actual free gauge field plus six free Dirac fields, their candidate
stress tensors and the observables that distinguish them within this
framework. That comparison has not been completed.

\textbf{Current identities and normalization.} The checks preserve the
unrescaled axial and vector divergence identities when other fields are
inserted in the same correlation function. They also show that the
density rescalings are positive and invertible. Applying the mass-fit
stability estimates to the target's actual masses bounds the
pseudoscalar rescaling factor's error by \(12/1489\), about \(0.806\%\);
the scalar factor's bound is \(3/98\). These are bounds on rescaling the
scalar and pseudoscalar densities, not new normalizations of the
currents. Matching the conventional Wilson currents and controlling
residual terms at finite lattice spacing remain open.

\textbf{Applicability to the constructed theory.} Separate acceptance
theorems take the actual objects supplied by
{\small \nolinkurl{interacting_six}} and show that they retain all
original properties together with the explicit correlation limit,
normalization bounds, leading coupling-reading slope and eventual
positivity of that reading. This checks that the estimates apply to the
constructed objects, rather than only to hypothetical objects satisfying
extra conditions. The proofs still rely on the submitted existence
theorem.

\subsection{One-loop flavor dependence and the meaning of the
coupling}\label{one-loop-flavor-dependence-and-the-meaning-of-the-coupling}

The beta function describes how a renormalized coupling changes with
scale. Its leading perturbative term is the \emph{one-loop} term. In QCD
that term depends on the number of quark flavors, denoted \(N_f\). The
source's trace calculations and its running-parameter formula have the
expected dependence:

\[
B_0(N_f)=11-\frac{2}{3}N_f,\qquad
b_0=\frac{B_0}{16\pi^2},\qquad
\operatorname{betaQ}(N_f)=12b_0.
\]

The last equality translates to the source's normalization; it should
not be read as a definition of a conventional physical beta function.
The
\href{https://github.com/yang-mills-group/QCD/blob/main/docs/flavor-running-tests.md}{flavor
and running development}, which also includes auxiliary reference
calculations, adds 136 distinct audited declarations across five suites,
27 Comparator targets and 26 specification controls.

There is also a nontrivial fermion calculation behind the checks. Lean
proofs establish a nonzero screening response of the free Wilson fermion
determinant in a fixed gauge background, and the multiplicity of that
response for repeated flavors. They include a probe on which the
response is nonzero. The benchmark uses identical masses for its
determinant calculation; the final existence theorem instead has six
positive distinct scheme masses. This supports one-loop flavor
dependence without completing the identification of the source's running
parameter with a physical renormalized coupling.

That distinction matters because the target uses a particular coupling
\emph{reading}, called \textbf{stiffness}. It uses second derivatives
near flat gauge links of an action defined as minus the logarithm of the
real part of the blocked density's scalar Grassmann coefficient
({\small \nolinkurl{Tower.lean}}). The construction keeps this reading
within a prescribed tolerance of its coupling parameter and proves its
leading logarithmic slope. It does not identify the parameter exactly
with a uniquely chosen physical coupling scheme.

Taking the scalar coefficient of a Grassmann expression and integrating
over its fermionic variables are different operations. Our checks make
this explicit: at the original lattice scale and a fixed bare gauge
parameter, the scalar-coefficient reading is unchanged when flavors or
masses change, whereas fermionic integration produces a gauge-field
weight containing the full product of fermion determinants. This finding
does not imply that all coarse-scale readings ignore flavor or that the
construction omits fermions.

One further limit concerns the way the proof schedules changes of
lattice scale. Its active condition, named {\small \nolinkurl{Earlier}},
requires \(N_f\leq8\). A recorded inspection of the Lean declaration
dependencies locates this condition and the later slope theorem among
the dependencies of the accepted result. This restricts the estimates
used here; it is not a general obstruction to a different construction.
The existence theorem verified in this review is specifically the
six-flavor theorem.

\subsection{Two-loop derivation
deferred}\label{two-loop-derivation-deferred}

The two-loop coefficient is the next perturbative coefficient in the
running of the coupling. The verified leading logarithmic slope does not
determine it. Our control examples show that different subleading
log-log terms can share that slope. Even fixing the first two asymptotic
coefficients and imposing monotonicity does not determine a differential
equation for the coupling.

\textbf{Deriving the physical two-loop coefficient from the constructed
theory is deferred to later work.} It requires identifying the physical
coupling and proving more precise control of its derivative or its
change between scales. The current source estimates and review checks do
not supply that derivation.

\subsection{How this relates to the informal proof
review}\label{how-this-relates-to-the-informal-proof-review}

The included
\href{https://github.com/yang-mills-group/QCD/blob/main/docs/informal-part-VI-review-r02.pdf}{informal
Part VI review} is an 18-page analytic audit dated November 15. It
expressly says that it did not establish the prose manuscript's Theorem
96.1 and did not examine the Part VI Lean code. It identifies unresolved
estimates for renormalization-group evolution with field insertions,
pressure, masses, stress-tensor mixing and preservation of a common
vacuum, while checking some local and conditional algebra.

Those unresolved items do not mean that the Lean verification failed.
Conversely, successful Lean verification does not retroactively verify
each estimate or normalization choice in the prose manuscript. The
formal target allows a different route to the stated result, and the
mass-fit example above shows why the two accounts must be compared
explicitly.

The additional proofs now establish the stated separated-point limits
and normalization bounds for actual common objects. Remaining work
includes identifying the physical operators and coupling, completing the
free gauge-and-quark comparison, and deriving the physical two-loop
coefficient. The all-pairs mass fit should be stated wherever the
six-pair informal prescription was previously used. The finite Wilson
diagnostics do not by themselves establish the identification of the
continuum fields with conventional Wilson-field observables.

The justified conclusion remains that \textbf{the substantial six-flavor
massive-QCD formal target passed its acceptance procedure, and the
additional checks passed within their stated scopes}. Complete physical
identification, verification of every informal estimate, arbitrary
prescribed masses, a mass gap and confinement would be stronger claims
than the evidence reviewed here supports.

\subsection{Verification record and complete
inventory}\label{verification-record-and-complete-inventory}

The submission contains 7,594 Lean proof modules and 7,597 manifest
inputs. Its frozen reference has 222 inputs; 220 shared inputs were
checked for agreement. The reference was built separately before the
submission-only source was added. The exported reference contained
320,175,328 bytes and matched the verifier's expected SHA-256 hash:

\begin{verbatim}
32f8e0d3322520e251c27a96c7123aab2d695f68bb15fa60190e460cbe063dba
\end{verbatim}

The export had the same hash after the solution was added. The full
Comparator build completed 16,250 jobs, exited successfully and printed
all three required acceptance markers. The allowed axioms were
{\small \nolinkurl{propext}}, {\small \nolinkurl{Classical.choice}} and
{\small \nolinkurl{Quot.sound}}. File and tool hashes were checked again
around completion.

The run lasted approximately 6 hours 58 minutes, from 2026-11-16
17:36:04 UTC to 2026-11-17 00:34:37 UTC. Completion was 7:34 p.m. on
November 16 in New York. The
\href{https://github.com/yang-mills-group/QCD/blob/main/evidence/supplied-target/status.json}{status},
\href{https://github.com/yang-mills-group/QCD/blob/main/evidence/supplied-target/comparator.json}{Comparator
receipt} and
\href{https://github.com/yang-mills-group/QCD/blob/main/evidence/supplied-target/comparator.log}{full
log} are retained.

The build used Lean 4.32.0
({\small \nolinkurl{8c9756b28d64dab099da31a4c09229a9e6a2ef35}}) and
mathlib revision
{\small \nolinkurl{81a5d257c8e410db227a6665ed08f64fea08e997}}. Tool
provenance is recorded with the evidence. The macOS Nanoda executable
came from the checked official Lean 4.35.0-rc2 platform archive; that
archive supplied the checker, not the Lean version used to compile the
submission.

{\def\LTcaptype{none} % do not increment counter
\begin{longtable}[]{@{}
  >{\raggedright\arraybackslash}p{(\linewidth - 2\tabcolsep) * \real{0.3000}}
  >{\raggedright\arraybackslash}p{(\linewidth - 2\tabcolsep) * \real{0.7000}}@{}}
\toprule\noalign{}
\begin{minipage}[b]{\linewidth}\raggedright
Check
\end{minipage} & \begin{minipage}[b]{\linewidth}\raggedright
Completed result
\end{minipage} \\
\midrule\noalign{}
\endhead
\bottomrule\noalign{}
\endlastfoot
Input consistency & 7,597 submission inputs and 222 frozen-reference
inputs checked against recorded hashes; 7,594 submitted Lean proof
modules \\
Preserved reference & Separately built reference export matched the
expected SHA-256 before and after adding the solution \\
Submitted QCD theorem & Comparator, Nanoda and Lean passed;
approximately 6 hours 58 minutes \\
Additional Lean review & 177 distinct audited declarations: 8 initial,
33 correspondence and 136 later additions; only the permitted standard
axioms \\
Additional challenge statements & 37 targets passed Comparator and both
kernels: 3 initial, 7 correspondence and 27 later targets; eight
principal Comparator runs including the submission \\
Mathematical diagnostics and controls & 125 passed: 44 mass-fit, 41
finite Wilson, 9 input-guard, 5 Comparator and 26 later specification
controls \\
Later specification controls & Within the 125 above: 5 unchanged cases
accepted and 21 deliberately altered specifications rejected with the
intended proof failures \\
Evidence packaging & Integrity checks passed locally and in a fresh
published checkout; 13 additional packaging regression tests passed \\
\end{longtable}
}

These totals cover the complete review archive, including helper lemmas
and auxiliary reference calculations beyond the QCD findings discussed
here. The 177 declaration audits are Lean audits; only the separately
listed 37 challenge statements received the additional Comparator and
two-kernel check. The 125 diagnostics include numerical tests, exact
arithmetic and deliberate-error controls; they are not 125 formal
theorems. Rejection of a modified specification shows that the existing
proof detects that change, not that the modified proposition is false.
The 13 packaging tests are additional to the 125.

The full submission was built once. Later suites compiled new review
modules using its recorded build products. The
\href{https://github.com/yang-mills-group/QCD/blob/main/docs/evidence.md}{evidence
index} and
\href{https://github.com/yang-mills-group/QCD/blob/main/docs/reproduce.md}{reproduction
instructions} locate the source hashes, configurations, logs, axiom
reports, controls and source-overlap comparison.

The run used macOS, with the lack of Linux process isolation accepted
for this review. The {\small \nolinkurl{fake-landrun}} compatibility
shim does not provide Linux Landlock/systemd containment. Source
consistency and deliberate-error controls do not reproduce that
isolation. Installed tools and build dependencies remain part of the
verification assumptions; no Linux run is claimed.

\subsection{Provenance and limits}\label{provenance-and-limits}

The user supplied the source locally. Hashes identify the recorded
files; they do not establish authorship. OpenAI's
\href{https://github.com/yang-mills-group/submissions/blob/main/docs/submission-dates.md}{submission chronology},
supplied in correspondence, dates delivery of the massive-QCD formalization
to 4 November 2026. This corrects the earlier provisional November 10 label.
Package timestamps and the absence of source history still leave the
proof-creation date unverified. Submission to the group is distinct from
public release. This provenance note was updated on 24 November 2026.

{\small \nolinkurl{provenance/submission.json}} records the source
identity, and
{\small \nolinkurl{provenance/massive-qcd-lean-source.tar.gz}} preserves
the received files. The repository also includes frozen definitions,
verification records, review scripts and original notices; compiled
artifacts and downloaded tools are excluded. Source paths refer to the
archive's {\small \nolinkurl{proof/}} directory unless otherwise
specified. Frozen definitions are directly available in
{\small \nolinkurl{provenance/supplied-comparator/frozen/reference/}}.

\textbf{Evidence snapshot.} Links to supporting records use repository
revision {\small \nolinkurl{81b8b764ec42}}. This editorial revision
changes no mathematical source or recorded verdict.

\end{document}
